\documentclass[11pt]{article}
\usepackage[margin=1in]{geometry}
\usepackage{amsmath,amssymb}
\usepackage{booktabs}
\usepackage{graphicx}
\usepackage{multirow}
\usepackage{xcolor}
\usepackage{hyperref}
\usepackage{natbib}
\usepackage{tabularx}
\usepackage{caption}

\title{Generation Depth Estimation: A Missing Layer in Synthetic Text Analysis}

\author{
  Anonymous Authors
}

\date{}

\begin{document}
\maketitle

% ============================================================
\begin{abstract}
% ============================================================
Web-scraped corpora increasingly contain LLM-generated text produced through recursive generation chains, where outputs of one model serve as inputs to the next generation step.
Current detection methods provide only binary real-vs-synthetic classification, obscuring the \emph{generation depth}---the number of recursive generation steps separating a text from its human-written seed.
We formalize \textbf{generation depth estimation} as an ordinal classification task and establish baselines using distributional features computed under a held-out reference language model.
On in-distribution data (single generator, single domain), a feature-based classifier achieves macro-F1 of 0.82 for 3-class depth estimation (gen-0/gen-1/gen-2+), and reference-model perplexity exhibits a robust monotonic decrease across depths (13.5 $\to$ 1.7).
However, systematic evaluation across 12 experimental conditions reveals severe limitations: cross-generator generalization drops to F1~0.52, reference-model depth correlation is family-specific ($\rho$=-0.68 same-family vs.\ $\rho$=-0.15 cross-family), and depth signal degrades with generator scale (F1 correlation with log-parameters: $r$=-0.85), with a 7B+ generator yielding near-chance performance (F1=0.43).
We identify text length as a major confound explaining 61--68\% of surface feature depth correlation, and show that finer-grained (5-class) depth estimation is near chance (F1=0.28).
These results establish generation depth estimation as a well-defined but challenging task, identify reference-model perplexity as the most promising depth proxy, and provide evidence that depth signal exhibits scale-dependent degradation consistent with model collapse theory.
We release benchmark data and evaluation protocols to support future work.
\end{abstract}

% ============================================================
\section{Introduction}
\label{sec:intro}
% ============================================================

Large language models are increasingly used in pipelines where one model's output becomes another's input---through automated content generation, data augmentation, or recursive summarization.
When such text enters web-scraped training corpora, it introduces contamination of varying severity.
A text generated directly from a human-written prompt (generation depth 1) may retain most properties of its human seed, while text at depth 4 or 5 has undergone compounding distributional drift through multiple recursive generation steps.

Current synthetic text detection methods~\citep{mitchell2023detectgpt,bao2024fastdetectgpt,hans2024binoculars} provide only binary real-vs-synthetic classification.
This forces a binary curation decision: discard all flagged content (losing potentially useful shallow-generation text) or retain everything (admitting deeply generated text with compounding quality degradation).
A more nuanced approach would estimate the \emph{generation depth} of each document, enabling risk-stratified data curation.

In this paper, we formalize \textbf{generation depth estimation} as an ordinal classification task and conduct a comprehensive empirical study of its feasibility using distributional features.
Our study spans 12 experimental conditions covering cross-generator generalization, reference model sensitivity, generator scale, prompt type variation, domain effects, ordinal granularity, edit robustness, and downstream utility.

Our main findings are:

\begin{enumerate}
\item \textbf{Task formulation}: We define generation depth estimation as a 3-class ordinal problem (gen-0/gen-1/gen-2+) and establish that this granularity represents a practical ceiling---5-class estimation is near chance (F1=0.28).

\item \textbf{Reference-model perplexity as depth proxy}: Median perplexity under a same-family reference model decreases monotonically with depth (13.5 $\to$ 1.7), providing a robust univariate depth signal. However, this signal is reference-model-specific: cross-family references yield weak correlations ($|\rho|<0.3$).

\item \textbf{Scale-dependent signal degradation}: Depth estimation F1 correlates strongly with generator scale ($r$=-0.85). For generators $\leq$3B, 3-class F1 ranges 0.68--0.78; at 8B, F1 drops to 0.43. This connects empirically to model collapse theory: larger models collapse slower, producing less distributional drift per recursive step.

\item \textbf{Comprehensive limitation analysis}: We document failure modes including cross-generator degradation (F1=0.52), length confounds (10-point F1 drop after control), prompt-type sensitivity, and fragility to post-generation editing.
\end{enumerate}

% ============================================================
\section{Related Work}
\label{sec:related}
% ============================================================

\paragraph{Synthetic text detection.}
Binary detection of machine-generated text has been approached through perplexity-based methods~\citep{gehrmann2019gltr,mitchell2023detectgpt,bao2024fastdetectgpt,hans2024binoculars}, trained classifiers~\citep{solaiman2019release,uchendu2020authorship}, watermarking~\citep{kirchenbauer2023watermark}, and character-level distributional signatures (LD-Score)~\citep{tulchinskii2024ldscore}.
All existing methods produce binary outputs; none attempt ordinal depth estimation.

\paragraph{Model collapse and recursive generation.}
Theoretical and empirical work on model collapse~\citep{alemohammad2024selfconsuming,shumailov2024curse,dohmatob2024model,gerstgrasser2024model} shows that recursive training on synthetic data leads to distributional convergence, with convergence rate inversely related to model capacity.
\citet{briesch2023attractor} prove that paraphrastic generation chains converge to 2-period attractor cycles, bounding useful depth signal to approximately 3 classes.
We extend this analysis empirically to non-paraphrastic chains.

\paragraph{Data quality and curation.}
Methods for training data curation~\citep{xie2023dsir,wettig2024qurating,wenzek2020ccnet} score documents along quality dimensions but do not account for generation provenance or depth.
Our work introduces generation depth as a complementary quality signal.

\paragraph{Vocabulary diversity in LLM outputs.}
Studies on vocabulary diversity in synthetic text~\citep{guo2024curious,padmakumar2024does,holtzman2020curious} document diversity reduction in LLM outputs.
We connect these findings to recursive generation, showing that diversity metrics (TTR, hapax ratio) decrease with generation depth but are substantially confounded by text length.

% ============================================================
\section{Task Definition}
\label{sec:task}
% ============================================================

\paragraph{Generation depth.}
Given a text $x$, its \emph{generation depth} $d(x) \in \{0, 1, 2, \ldots\}$ is defined as the number of LLM generation steps separating it from a human-written seed:
\begin{itemize}
  \item $d=0$: human-written text (the original seed).
  \item $d=1$: text generated by an LLM from a human-written prompt/seed.
  \item $d=k$: text produced by applying LLM generation $k$ times recursively.
\end{itemize}

\paragraph{Depth estimation task.}
Given a text $x$ with unknown provenance, predict $\hat{d}(x) \in \{$gen-0, gen-1, gen-2+$\}$.
We adopt 3-class ordinal grouping based on theoretical and empirical evidence that finer-grained estimation is infeasible with distributional features (Section~\ref{sec:ordinal}).

\paragraph{Generation chain construction.}
We construct generation chains via iterative continuation: given seed text $s$, we produce $x^{(1)} = \text{LLM}(s)$, $x^{(2)} = \text{LLM}(x^{(1)})$, and so on up to depth 4.
We explore continuation, elaboration, and inspired-rewriting prompt types (Section~\ref{sec:prompt_type}).
Chains may use a single generator (Section~\ref{sec:indist}--\ref{sec:scale}) or different generators per step (Section~\ref{sec:cross_model}).

\paragraph{Evaluation protocol.}
Primary metric: macro-averaged F1 across the 3 depth classes.
We report 95\% bootstrap confidence intervals and use leave-one-generator-out (LOGO) evaluation for cross-generator generalization.

% ============================================================
\section{Method}
\label{sec:method}
% ============================================================

\subsection{Distributional Features}
\label{sec:features}

We extract 7 features from each text, divided into two categories:

\paragraph{Surface features} computed directly from the text:
\begin{itemize}
  \item \textbf{TTR}: Type-token ratio (unique tokens / total tokens).
  \item \textbf{Hapax ratio}: Fraction of tokens appearing exactly once.
\end{itemize}

\paragraph{Reference-model features} computed by evaluating the text under a held-out reference language model $M_{\text{ref}}$ (default: Pythia-410M), collecting per-token negative log-likelihoods $\{-\log p_{M_{\text{ref}}}(t_i \mid t_{<i})\}_{i=1}^{n}$:
\begin{itemize}
  \item \textbf{entropy\_mean}: Mean NLL (proxy for reference-model perplexity).
  \item \textbf{entropy\_std}: Standard deviation of per-token NLL.
  \item \textbf{entropy\_max}: Maximum per-token NLL.
  \item \textbf{entropy\_min}: Minimum per-token NLL.
  \item \textbf{tail\_decay}: Rate of exponential decay in the upper tail of the NLL distribution.
\end{itemize}

These features are standard NLL summary statistics; the novelty lies in applying them to ordinal depth estimation rather than binary detection.

\subsection{Classifier}
We train an XGBoost gradient-boosted tree classifier on the 7-dimensional feature vector.
Hyperparameters are selected via 5-fold cross-validation on training data.
We also evaluate a QLoRA-finetuned Qwen3.5-0.8B as an end-to-end baseline (Section~\ref{sec:e2e}).

% ============================================================
\section{Experiments}
\label{sec:experiments}
% ============================================================

We organize experiments around four research questions:
\textbf{RQ1}: Can distributional features estimate generation depth? (Sections~\ref{sec:indist}--\ref{sec:length})
\textbf{RQ2}: How does depth signal scale with generator capacity? (Section~\ref{sec:scale})
\textbf{RQ3}: What are the generalization boundaries? (Sections~\ref{sec:cross_gen}--\ref{sec:cross_model})
\textbf{RQ4}: Is depth estimation useful for downstream curation? (Section~\ref{sec:downstream})

% --------------------------------------------------
\subsection{In-Distribution Baseline}
\label{sec:indist}

We first establish baseline performance using Pythia-2.8B as generator and Pythia-410M as reference model on Wikipedia-domain continuation chains.

\begin{table}[h]
\centering
\caption{Feature ablation on in-distribution data (Pythia-2.8B generator, Pythia-410M reference, Wikipedia continuation chains).}
\label{tab:ablation}
\begin{tabular}{lccccc}
\toprule
Feature Set & Macro-F1 & 95\% CI & Gen-0 & Gen-1 & Gen-2+ \\
\midrule
Full (7 features) & 0.820 & [0.807, 0.834] & 0.952 & 0.776 & 0.795 \\
Surface only & 0.763 & [0.747, 0.774] & 0.920 & 0.710 & 0.658 \\
Reference only & 0.719 & [0.706, 0.733] & 0.895 & 0.640 & 0.621 \\
\bottomrule
\end{tabular}
\end{table}

In-distribution, the full feature set achieves macro-F1 of 0.820 (Table~\ref{tab:ablation}).
Surface features alone yield 0.763 and reference features alone 0.719, indicating complementary signal.
Median reference-model perplexity decreases monotonically across depths: 13.5 (gen-0), 4.96 (gen-1), 3.00 (gen-2), 2.07 (gen-3), 1.72 (gen-4).

Feature-level analysis reveals two behavioral groups (Table~\ref{tab:h1h2}).
Three features show monotonic depth correlation ($R^2 > 0.8$): entropy\_mean ($\rho$=-0.68), TTR ($\rho$=-0.31), and hapax\_ratio ($\rho$=-0.28).
Four features show negligible depth correlation ($R^2 < 0.1$): entropy\_std, entropy\_max, entropy\_min, and tail\_decay.
We label these drifting (H2) and stationary (H1) groups respectively.

\begin{table}[h]
\centering
\caption{Feature classification by depth behavior. H2: monotonic depth drift. H1: near-stationary.}
\label{tab:h1h2}
\begin{tabular}{lcccc}
\toprule
Feature & Source & Depth Slope & $R^2$ & Class \\
\midrule
entropy\_mean & reference & $-0.042$ & 0.91 & H2 (drift) \\
TTR & surface & $-0.031$ & 0.87 & H2 (drift) \\
hapax\_ratio & surface & $-0.028$ & 0.83 & H2 (drift) \\
\midrule
entropy\_std & reference & $+0.003$ & 0.04 & H1 (stationary) \\
entropy\_max & reference & $-0.005$ & 0.07 & H1 (stationary) \\
entropy\_min & reference & $+0.002$ & 0.02 & H1 (stationary) \\
tail\_decay & reference & $+0.004$ & 0.05 & H1 (stationary) \\
\bottomrule
\end{tabular}
\end{table}

However, as subsequent sections will show, these in-distribution results are misleadingly optimistic due to same-family reference model alignment, length confounds, and single-generator specificity.

% --------------------------------------------------
\subsection{Text Length as Confound}
\label{sec:length}

Generated texts exhibit a length-depth correlation ($r$=0.42): deeper generations tend to be shorter.
Since surface diversity metrics are mechanically length-dependent, this confound must be addressed.

\begin{table}[h]
\centering
\caption{Classification performance before and after length control.}
\label{tab:length}
\begin{tabular}{lcccc}
\toprule
Condition & Full F1 & Surface F1 & Reference F1 & Gap \\
\midrule
Original & 0.78 & 0.76 & 0.72 & 0.02 \\
Length-normalized features & 0.68 & 0.63 & 0.66 & 0.02 \\
Length-matched subsampling & 0.70 & 0.65 & 0.67 & 0.03 \\
\bottomrule
\end{tabular}
\end{table}

Length normalization reduces full-feature F1 from 0.78 to 0.68---a 10-point drop (Table~\ref{tab:length}).
The impact is concentrated in surface features: TTR depth correlation drops from $\rho$=-0.31 to $\rho$=-0.12 (61\% explained by length), and hapax ratio from $\rho$=-0.28 to $\rho$=-0.09 (68\% explained by length) as shown in Table~\ref{tab:length_features}.
Reference-model entropy\_mean is less affected (19\% explained by length, $\rho$: $-0.68 \to -0.55$), confirming it carries genuine depth signal beyond the length confound.

\begin{table}[h]
\centering
\caption{Feature depth correlation before and after length normalization.}
\label{tab:length_features}
\begin{tabular}{lccc}
\toprule
Feature & Original $\rho$ & Length-norm.\ $\rho$ & \% Explained by Length \\
\midrule
TTR & $-0.31$ & $-0.12$ & 61\% \\
hapax\_ratio & $-0.28$ & $-0.09$ & 68\% \\
entropy\_mean & $-0.68$ & $-0.55$ & 19\% \\
entropy\_std & $+0.03$ & $+0.02$ & 33\% \\
tail\_decay & $+0.04$ & $+0.03$ & 25\% \\
\bottomrule
\end{tabular}
\end{table}

After controlling for length, the advantage of the full feature set over surface-only features nearly vanishes (0.68 vs 0.63 = 0.05, down from 0.06), indicating that reference features' previous advantage was partly due to being less length-confounded rather than capturing fundamentally different depth information.

% --------------------------------------------------
\subsection{Cross-Generator Generalization}
\label{sec:cross_gen}

We evaluate generalization using leave-one-generator-out (LOGO) across 5 generators from 3 model families: Pythia-2.8B, Pythia-1B, OLMo-2-1B, Qwen3-1.7B, and Qwen3-0.6B.

\begin{table}[h]
\centering
\caption{Leave-one-generator-out macro-F1. Mean across generators: 0.52.}
\label{tab:logo}
\begin{tabular}{llcccc}
\toprule
Held-Out Generator & Family & Macro-F1 & Gen-0 & Gen-1 & Gen-2+ \\
\midrule
Pythia-2.8B & Pythia & 0.61 & 0.82 & 0.51 & 0.50 \\
Pythia-1B & Pythia & 0.64 & 0.85 & 0.54 & 0.53 \\
OLMo-2-1B & OLMo & 0.48 & 0.72 & 0.38 & 0.34 \\
Qwen3-1.7B & Qwen & 0.44 & 0.68 & 0.35 & 0.29 \\
Qwen3-0.6B & Qwen & 0.43 & 0.67 & 0.33 & 0.29 \\
\midrule
Mean (all) & --- & 0.52 & 0.75 & 0.42 & 0.39 \\
\bottomrule
\end{tabular}
\end{table}

LOGO macro-F1 averages 0.52 (Table~\ref{tab:logo}), substantially below the in-distribution 0.82.
Within-family generalization (Pythia $\to$ Pythia) averages 0.625, while cross-family generalization (e.g., Pythia+Qwen $\to$ OLMo) averages 0.45.
Gen-0 detection remains passable (mean 0.75) because human-written text is universally distinct, but gen-1 vs gen-2+ discrimination---the core depth signal---collapses across families.

\begin{table}[h]
\centering
\caption{Leave-one-family-out evaluation.}
\label{tab:lofo}
\begin{tabular}{lcc}
\toprule
Held-Out Family & Train Generators & Macro-F1 \\
\midrule
Pythia & OLMo+Qwen & 0.41 \\
OLMo & Pythia+Qwen & 0.54 \\
Qwen & Pythia+OLMo & 0.56 \\
\bottomrule
\end{tabular}
\end{table}

Leave-one-family-out results (Table~\ref{tab:lofo}) confirm family-level dependence: holding out Pythia (the family matching our reference model) drops F1 to 0.41, indicating the classifier critically relies on same-family reference model alignment.

% --------------------------------------------------
\subsection{Reference Model Sensitivity}
\label{sec:ref_model}

We test whether the depth signal generalizes across reference models by computing features using OLMo-2-1B and Qwen3-0.6B in addition to the default Pythia-410M.

\begin{table}[h]
\centering
\caption{Entropy\_mean depth correlation by reference model. The strong H2 drift observed with Pythia-410M is reference-family-specific.}
\label{tab:ref_corr}
\begin{tabular}{llccc}
\toprule
Reference Model & Family & entropy\_mean $\rho$ & $p$-value & H2 Status \\
\midrule
Pythia-410M & Pythia & $-0.68$ & $<$0.001 & Confirmed \\
OLMo-2-1B & OLMo & $-0.22$ & 0.031 & Weak ($|\rho|<0.3$) \\
Qwen3-0.6B & Qwen & $-0.15$ & 0.118 & Not confirmed \\
\bottomrule
\end{tabular}
\end{table}

The strong entropy\_mean depth correlation ($\rho$=-0.68) observed with the Pythia-410M reference largely disappears with non-Pythia references: $\rho$=-0.22 (OLMo) and $\rho$=-0.15 (Qwen, non-significant; Table~\ref{tab:ref_corr}).
This occurs because the Pythia-410M reference shares architecture and training data with the Pythia generators, producing calibrated perplexity estimates that systematically track depth.
Cross-family references assign noisier, less calibrated perplexities.

\begin{table}[h]
\centering
\caption{Classification performance by reference model.}
\label{tab:ref_class}
\begin{tabular}{lcccc}
\toprule
Reference Model & Full F1 & Surface F1 & Ref-only F1 & Full--Surface Gap \\
\midrule
Pythia-410M & 0.82 & 0.76 & 0.72 & 0.06 \\
OLMo-2-1B & 0.63 & 0.61 & 0.48 & 0.02 \\
Qwen3-0.6B & 0.58 & 0.57 & 0.41 & 0.01 \\
\bottomrule
\end{tabular}
\end{table}

Classification F1 drops sharply with cross-family references (Table~\ref{tab:ref_class}): from 0.82 (Pythia) to 0.63 (OLMo) and 0.58 (Qwen).
Reference-only F1 drops to 0.48 and 0.41 respectively, while surface-only F1 remains stable at 0.57--0.61 (as expected, since surface features are reference-model-independent).

Examining feature stability across references (Table~\ref{tab:stability}), we find that distributional shape features (entropy\_std, entropy\_max, tail\_decay) remain H1-stationary regardless of reference model---this is genuinely robust.
However, the drifting features reduce to surface-only metrics (TTR, hapax), which Section~\ref{sec:length} showed are largely length artifacts.

\begin{table}[h]
\centering
\caption{Feature behavioral classification stability across reference models.}
\label{tab:stability}
\begin{tabular}{lcccc}
\toprule
Feature & Pythia Class & OLMo Class & Qwen Class & Stable? \\
\midrule
entropy\_mean & H2 & weak/NS & NS & No \\
entropy\_std & H1 & H1 & H1 & Yes \\
entropy\_max & H1 & H1 & H1 & Yes \\
tail\_decay & H1 & H1 & H1 & Yes \\
TTR & H2 & H2 & H2 & Yes \\
hapax\_ratio & H2 & H2 & weak H2 & Marginal \\
\bottomrule
\end{tabular}
\end{table}

% --------------------------------------------------
\subsection{Generator Scale Analysis}
\label{sec:scale}

We evaluate depth estimation across generator scales from 0.6B to 8B parameters (Table~\ref{tab:scale}).

\begin{table}[h]
\centering
\caption{Depth classification by generator scale. Depth signal degrades with model capacity ($r$=-0.85 between log-params and F1).}
\label{tab:scale}
\begin{tabular}{lcccccc}
\toprule
Generator & Params & F1 & Gen-0 & Gen-1 & Gen-2+ & $\rho_{\text{ent}}$ \\
\midrule
Qwen3-0.6B & 0.6B & 0.74 & 0.91 & 0.66 & 0.65 & $-0.58$ \\
Pythia-1B & 1B & 0.71 & 0.89 & 0.62 & 0.62 & $-0.55$ \\
OLMo-2-1B & 1.5B & 0.69 & 0.87 & 0.60 & 0.60 & $-0.51$ \\
Qwen3-1.7B & 1.7B & 0.68 & 0.86 & 0.58 & 0.60 & $-0.48$ \\
Pythia-2.8B & 2.8B & 0.78 & 0.93 & 0.72 & 0.69 & $-0.68$ \\
Qwen3-8B & 8B & 0.43 & 0.62 & 0.35 & 0.32 & $-0.19$ \\
\bottomrule
\end{tabular}
\end{table}

The 8B generator yields near-chance depth estimation (F1=0.43) with severely degraded gen-1 and gen-2+ discrimination (0.35 and 0.32).
The correlation between log-parameters and macro-F1 is $r$=-0.85 (excluding the Pythia-2.8B outlier, which benefits from same-family reference alignment).

The mechanism is revealed by perplexity dynamic range analysis (Table~\ref{tab:ppl_range}): the reference model's ability to distinguish depths relies on perplexity separation between depths, which compresses as generator fluency increases.
For Pythia-2.8B, perplexity ranges from 13.5 (gen-0) to 1.7 (gen-4), a dynamic range of 11.8.
For Qwen3-8B, this compresses to 8.2--4.1, a dynamic range of only 4.1---the larger model produces text that is already close to the reference model's expectations at every depth.

\begin{table}[h]
\centering
\caption{Reference-model perplexity dynamic range by generator scale. Larger generators compress the perplexity range available for depth discrimination.}
\label{tab:ppl_range}
\begin{tabular}{lcccc}
\toprule
Generator & PPL (depth-0) & PPL (depth-4) & Dynamic Range & PPL Ratio \\
\midrule
Pythia-2.8B & 13.5 & 1.7 & 11.8 & 7.9$\times$ \\
Pythia-1B & 12.8 & 2.1 & 10.7 & 6.1$\times$ \\
OLMo-2-1B & 11.4 & 2.5 & 8.9 & 4.6$\times$ \\
Qwen3-1.7B & 10.8 & 2.8 & 8.0 & 3.9$\times$ \\
Qwen3-8B & 8.2 & 4.1 & 4.1 & 2.0$\times$ \\
\bottomrule
\end{tabular}
\end{table}

This result connects to model collapse theory~\citep{gerstgrasser2024model}: larger models collapse slower under recursive generation, producing less distributional drift per step.
Consequently, depth estimation with distributional features faces a fundamental scaling barrier---the models most likely to produce recursive content in practice (large, capable models) are precisely those for which depth signal is weakest.

% --------------------------------------------------
\subsection{Prompt Type and Chain Structure}
\label{sec:prompt_type}

We evaluate three prompt types: continuation (direct text continuation), elaboration (topic expansion), and inspired rewriting (semantic departure from source).

\begin{table}[h]
\centering
\caption{Depth classification by prompt type.}
\label{tab:prompt}
\begin{tabular}{lcccc}
\toprule
Prompt Type & Macro-F1 & $\rho_{\text{ent}}$ & $\rho_{\text{TTR}}$ & Gen-1 vs Gen-2+ F1 \\
\midrule
Continuation & 0.78 & $-0.65$ & $-0.29$ & 0.71 \\
Elaboration & 0.55 & $-0.38$ & $-0.16$ & 0.44 \\
Inspired rewriting & 0.47 & $-0.18$ & $-0.08$ & 0.36 \\
\bottomrule
\end{tabular}
\end{table}

Depth signal varies dramatically with prompt type (Table~\ref{tab:prompt}).
Continuation chains preserve the strongest signal (F1=0.78) because each step builds directly on the previous output, creating cumulative distributional drift.
Elaboration reduces this to F1=0.55, and inspired rewriting yields near-chance F1=0.47 because the semantic departure partially resets distributional state at each step.

Cross-prompt transfer is poor: training on one prompt type and testing on another yields average F1=0.42 (Table~\ref{tab:prompt_transfer}), indicating that depth features are prompt-type-specific.

\begin{table}[h]
\centering
\caption{Cross-prompt type transfer (train $\to$ test).}
\label{tab:prompt_transfer}
\begin{tabular}{lccc}
\toprule
Train $\to$ Test & Elaboration & Rewriting & Continuation \\
\midrule
Continuation $\to$ & 0.46 & 0.39 & --- \\
Elaboration $\to$ & --- & 0.42 & 0.51 \\
Rewriting $\to$ & 0.40 & --- & 0.44 \\
\bottomrule
\end{tabular}
\end{table}

\paragraph{Paraphrastic vs.\ non-paraphrastic chains.}
Comparing paraphrastic and continuation chains (Table~\ref{tab:para}) confirms attractor theory predictions: paraphrastic chains converge by depth 2 (gen-2 vs gen-3 F1=0.34), while continuation chains preserve signal slightly longer but also saturate by depth 3.

\begin{table}[h]
\centering
\caption{Paraphrastic vs.\ continuation chain depth discrimination.}
\label{tab:para}
\begin{tabular}{lcccccc}
\toprule
Chain Type & 3-class F1 & Gen-0 & Gen-1 & Gen-2+ & $\rho_{\text{ent}}$ \\
\midrule
Paraphrastic & 0.52 & 0.78 & 0.45 & 0.33 & $-0.31$ \\
Continuation & 0.78 & 0.93 & 0.72 & 0.69 & $-0.65$ \\
\bottomrule
\end{tabular}
\end{table}

% --------------------------------------------------
\subsection{Cross-Model Chains}
\label{sec:cross_model}

Real-world contamination may involve different generators at each depth step.
We evaluate cross-model chains where the generator changes randomly at each step (Table~\ref{tab:crossmodel}).

\begin{table}[h]
\centering
\caption{Single-model vs.\ cross-model chain depth estimation.}
\label{tab:crossmodel}
\begin{tabular}{lcccc}
\toprule
Chain Type & Macro-F1 & Gen-0 & Gen-1 & Gen-2+ \\
\midrule
Single model & 0.78 & 0.93 & 0.72 & 0.69 \\
Cross-model (random) & 0.39 & 0.65 & 0.28 & 0.24 \\
Cross-model (increasing size) & 0.41 & 0.66 & 0.30 & 0.27 \\
Cross-model (decreasing size) & 0.38 & 0.63 & 0.27 & 0.24 \\
\bottomrule
\end{tabular}
\end{table}

Cross-model chains destroy depth signal (F1=0.39), with gen-1 and gen-2+ near chance (0.28 and 0.24).
Generator identity is a much stronger signal than depth: a gen-1 text from Qwen differs more from gen-1 text from Pythia than gen-1 differs from gen-2 within the same generator.
This represents the most ecologically challenging result, as real web contamination involves diverse generators.

% --------------------------------------------------
\subsection{Multi-Domain Evaluation}
\label{sec:domain}

We expand to 4 domains with 500 seeds each (Table~\ref{tab:domain}).

\begin{table}[h]
\centering
\caption{Depth classification by source domain.}
\label{tab:domain}
\begin{tabular}{lccccc}
\toprule
Domain & $n$ Seeds & Macro-F1 & Gen-0 & Gen-1 & Gen-2+ \\
\midrule
Wikipedia & 500 & 0.72 & 0.89 & 0.63 & 0.64 \\
News & 500 & 0.68 & 0.86 & 0.59 & 0.59 \\
Scientific & 500 & 0.55 & 0.78 & 0.44 & 0.43 \\
Creative & 500 & 0.48 & 0.71 & 0.37 & 0.36 \\
\midrule
All pooled & 2000 & 0.61 & 0.81 & 0.51 & 0.51 \\
\bottomrule
\end{tabular}
\end{table}

Domain substantially moderates depth signal ($\sigma_{F1}$=0.11).
Factual domains (Wikipedia, news) show passable performance, while creative writing yields near-chance F1=0.48.
Cross-domain transfer averages F1=0.49, adding domain to the list of moderating factors (alongside prompt type and generator family) that limit generalizability.

% --------------------------------------------------
\subsection{Ordinal Granularity}
\label{sec:ordinal}

We test 5-class ordinal classification (gen-0 through gen-4) to establish the granularity ceiling (Table~\ref{tab:ordinal}).

\begin{table}[h]
\centering
\caption{Five-class ordinal depth classification.}
\label{tab:ordinal}
\begin{tabular}{lcccc}
\toprule
Class & Precision & Recall & F1 & $n$ \\
\midrule
Gen-0 & 0.88 & 0.82 & 0.85 & 995 \\
Gen-1 & 0.39 & 0.45 & 0.42 & 993 \\
Gen-2 & 0.16 & 0.20 & 0.18 & 912 \\
Gen-3 & 0.11 & 0.13 & 0.12 & 856 \\
Gen-4 & 0.10 & 0.12 & 0.11 & 816 \\
\midrule
\multicolumn{3}{l}{5-class macro-F1: 0.28} & \multicolumn{2}{l}{MAE: 1.4} \\
\bottomrule
\end{tabular}
\end{table}

Five-class macro-F1 of 0.28 is barely above chance (0.20).
Gen-2/3/4 collapse into an indistinguishable cluster (F1: 0.18/0.12/0.11), with 72\% of misclassifications between adjacent classes.
This validates the 3-class design: ordinal depth beyond gen-0/gen-1/gen-2+ is not discriminable with distributional features, consistent with attractor convergence theory extended to non-paraphrastic generation.

% --------------------------------------------------
\subsection{Robustness to Human Editing}
\label{sec:robustness}

\begin{table}[h]
\centering
\caption{Robustness to simulated post-generation editing.}
\label{tab:edits}
\begin{tabular}{lccc}
\toprule
Edit Type & Edit Rate & Macro-F1 & $\Delta$F1 \\
\midrule
None & 0\% & 0.78 & --- \\
Random replace & 5\% & 0.68 & $-0.10$ \\
Random replace & 10\% & 0.58 & $-0.20$ \\
Random replace & 20\% & 0.44 & $-0.34$ \\
LLM light edit & $\sim$8\% & 0.53 & $-0.25$ \\
\bottomrule
\end{tabular}
\end{table}

Depth features are fragile to post-generation modification (Table~\ref{tab:edits}).
Even 5\% random word replacement reduces F1 by 0.10, and LLM-simulated editing reduces F1 by 0.25.
LLM editing is more damaging than random replacement because it systematically smooths the distributional artifacts that depth estimation relies on.
This limits practical applicability to detection of fully automated, unedited generation chains.

% --------------------------------------------------
\subsection{Additional Features and End-to-End Baseline}
\label{sec:e2e}

\paragraph{LD-Score as depth feature.}
Adding LD-Score~\citep{tulchinskii2024ldscore} (a character-level distributional signature) provides negligible improvement: depth correlation $\rho$=-0.25, F1 increase from 0.780 to 0.785.
LD-Score shows weak depth sensitivity but insufficient signal to meaningfully contribute beyond existing features.

\begin{table}[h]
\centering
\caption{Handcrafted features vs.\ end-to-end fine-tuned classifier.}
\label{tab:e2e}
\begin{tabular}{lccc}
\toprule
Method & In-Dist.\ F1 & Cross-Gen LOGO F1 & GPU Hours \\
\midrule
Handcrafted + XGBoost & 0.78 & 0.52 & 0.1 \\
Qwen3.5-0.8B QLoRA & 0.83 & 0.55 & 15 \\
Qwen3.5-0.8B (frozen) & 0.80 & 0.53 & 8 \\
\bottomrule
\end{tabular}
\end{table}

\paragraph{End-to-end classifier.}
A QLoRA-finetuned Qwen3.5-0.8B achieves F1=0.83 in-distribution, slightly above the handcrafted 0.78 (Table~\ref{tab:e2e}).
However, it shares the same cross-generator failure (LOGO F1=0.55 vs.\ 0.52), suggesting that the generalization barrier is fundamental to the depth estimation task, not an artifact of our feature engineering.

% --------------------------------------------------
\subsection{Downstream Utility}
\label{sec:downstream}

We fine-tune OLMo-1B on depth-stratified data (5 seeds, bootstrap CIs) to test whether depth-aware filtering improves downstream model quality compared to binary filtering (Table~\ref{tab:downstream}).

\begin{table}[h]
\centering
\caption{Downstream fine-tuning results (OLMo-1B, 5 seeds).}
\label{tab:downstream}
\begin{tabular}{lcccc}
\toprule
Filtering Strategy & Mean PPL & Std & 95\% CI & Improvement \\
\midrule
Unfiltered & 14.8 & 0.42 & [14.2, 15.4] & --- \\
Binary (FastDetectGPT) & 14.5 & 0.51 & [13.8, 15.2] & 1.2\% [$-$1.1, 3.5] \\
Binary (LD-Score) & 14.6 & 0.48 & [14.0, 15.2] & 0.9\% [$-$1.4, 3.2] \\
Depth-aware (keep gen-0/1) & 14.2 & 0.39 & [13.7, 14.7] & 1.8\% [$-$0.5, 4.1] \\
Depth-aware (sample-matched) & 14.2 & 0.44 & [13.6, 14.8] & 1.5\% [$-$0.8, 3.8] \\
\bottomrule
\end{tabular}
\end{table}

Depth-aware filtering shows a directional advantage (1.8\% vs.\ 1.2\% improvement over unfiltered) but confidence intervals overlap substantially (depth vs.\ binary $p$=0.31).
The improvement is not statistically significant at this sample size.

Predicted depth does correlate weakly with quality degradation ($\rho$=0.28, 95\% CI [0.11, 0.45]).
Perplexity increases monotonically with depth: 12.1 (gen-0) $\to$ 13.4 (gen-1) $\to$ 14.1 (gen-2) $\to$ 14.5 (gen-3) $\to$ 14.9 (gen-4), confirming that deeper generation produces lower-quality training data.
However, this descriptive finding does not translate into significant downstream filtering advantage with current methods.

% ============================================================
\section{Discussion}
\label{sec:discussion}
% ============================================================

\paragraph{The task is well-defined but harder than expected.}
Generation depth estimation is conceptually sound: depth correlates with quality degradation, and ordinal estimation enables more nuanced curation than binary detection.
However, our comprehensive evaluation reveals that current distributional features fall short of reliable depth estimation under realistic conditions.
The 3-class problem is feasible in controlled settings (single generator, single domain, continuation chains) but degrades along every generalization axis we tested.

\paragraph{Reference-model perplexity: strong but specific.}
The monotonic perplexity decrease across depths (13.5 $\to$ 1.7) is the most robust univariate depth signal.
However, this signal depends critically on reference-generator family alignment ($\rho$=-0.68 same-family vs.\ $\rho$=-0.15 cross-family).
This aligns with an interpretation based on distributional calibration: a same-family reference model has calibrated predictions for the generator's output distribution, making perplexity sensitive to the subtle distributional shifts across depths.
A cross-family reference assigns noisier perplexities that wash out the depth signal.

\paragraph{Scale-dependent depth estimation.}
The strong negative correlation between generator scale and depth discriminability ($r$=-0.85) is our most concerning finding for practical deployment.
Larger models produce more fluent text with smaller distributional drift per recursive step, compressing the perplexity dynamic range available for depth discrimination (11.8 for 2.8B vs.\ 4.1 for 8B).
This connects to theoretical results on scale-dependent collapse rates~\citep{gerstgrasser2024model}: if larger models converge slower, then depth-correlated distributional drift is correspondingly smaller at each step.
As the most capable models are also those most likely to be used in recursive generation pipelines, this scaling barrier is a fundamental challenge.

\paragraph{Confounds and moderators.}
We identify text length, prompt type, generator family, source domain, and post-generation editing as significant moderators.
The length confound is particularly concerning: 61--68\% of surface feature depth correlation is attributable to length variation.
After controlling for these confounds, the residual depth signal, while real, is substantially weaker than in-distribution baselines suggest.

\paragraph{Limitations.}
Our study is limited to open-weight models $\leq$8B parameters; frontier models (70B+) may exhibit even weaker depth signal.
The benchmark uses constructed chains rather than naturally occurring recursive generation, which may differ in distribution.
All reference models used are autoregressive transformers; alternative architectures may provide complementary signals.

% ============================================================
\section{Conclusion}
\label{sec:conclusion}
% ============================================================

We formalize generation depth estimation as a new task in synthetic text analysis and provide a comprehensive benchmark spanning 12 experimental conditions.
Our study identifies three key findings: (1)~three-class depth estimation (gen-0/gen-1/gen-2+) represents the practical granularity ceiling; (2)~reference-model perplexity is the strongest depth proxy but requires generator-family-aligned reference models; and (3)~depth signal exhibits scale-dependent degradation, with 8B generators already yielding near-chance estimation.
These results position generation depth estimation as a meaningful but challenging open problem, and the scale-dependent signal degradation connects empirically to model collapse theory.
We release our benchmark dataset and evaluation protocols to support future work on generator-invariant depth features that may overcome the limitations documented here.

% ============================================================
% References
% ============================================================

\bibliographystyle{plainnat}

\begin{thebibliography}{30}

\bibitem[Alemohammad et~al.(2024)]{alemohammad2024selfconsuming}
S.~Alemohammad, J.~Casco-Rodriguez, L.~Luber, A.~G. Humair, S.~Babaei, D.~LeJeune, A.~Sivaraman, and R.~G. Baraniuk.
\newblock Self-consuming generative models go {MAD}.
\newblock In \emph{ICLR}, 2024.

\bibitem[Bao et~al.(2024)]{bao2024fastdetectgpt}
G.~Bao, Y.~Zhao, Z.~Teng, L.~Yang, and Y.~Zhang.
\newblock {Fast-DetectGPT}: Efficient zero-shot detection of machine-generated text via conditional probability curvature.
\newblock In \emph{ICLR}, 2024.

\bibitem[Briesch et~al.(2023)]{briesch2023attractor}
M.~Briesch, D.~Sobania, and F.~Rothlauf.
\newblock Large language models suffer from their own output: An analysis of the self-consuming training loop.
\newblock \emph{arXiv preprint arXiv:2311.16822}, 2023.

\bibitem[Dohmatob et~al.(2024)]{dohmatob2024model}
E.~Dohmatob, P.~Feng, Q.~Yang, F.~Charton, and J.~Kempe.
\newblock A tale of tails: Model collapse as a change of scaling laws.
\newblock In \emph{ICML}, 2024.

\bibitem[Gehrmann et~al.(2019)]{gehrmann2019gltr}
S.~Gehrmann, H.~Strobelt, and A.~M. Rush.
\newblock {GLTR}: Statistical detection and visualization of generated text.
\newblock In \emph{ACL}, 2019.

\bibitem[Gerstgrasser et~al.(2024)]{gerstgrasser2024model}
M.~Gerstgrasser, R.~Schaeffer, A.~Dey, R.~Rafailov, H.~Tombai, Y.~Nishi, S.~Levine, and K.~Koyejo.
\newblock Is model collapse inevitable? {Breaking} the curse of recursion by accumulating real and synthetic data.
\newblock In \emph{NeurIPS}, 2024.

\bibitem[Guo et~al.(2024)]{guo2024curious}
B.~Guo, X.~Zhang, Z.~Wang, M.~Jiang, J.~Nie, Y.~Ding, J.~Yue, and Y.~Wu.
\newblock How close is {ChatGPT} to human experts? {Comparison} corpus, evaluation, and detection.
\newblock \emph{arXiv preprint arXiv:2301.07597}, 2024.

\bibitem[Hans et~al.(2024)]{hans2024binoculars}
A.~Hans, A.~Schwarzschild, V.~Cherepanova, H.~Kazemi, A.~Saha, M.~Goldblum, J.~Geiping, and T.~Goldstein.
\newblock Spotting {LLMs} with {Binoculars}: Zero-shot detection of machine-generated text.
\newblock In \emph{ICML}, 2024.

\bibitem[Hendrycks and Gimpel(2017)]{hendrycks2017baseline}
D.~Hendrycks and K.~Gimpel.
\newblock A baseline for detecting misclassified and out-of-distribution examples in neural networks.
\newblock In \emph{ICLR}, 2017.

\bibitem[Holtzman et~al.(2020)]{holtzman2020curious}
A.~Holtzman, J.~Buys, L.~Du, M.~Forbes, and Y.~Choi.
\newblock The curious case of neural text degeneration.
\newblock In \emph{ICLR}, 2020.

\bibitem[Kirchenbauer et~al.(2023)]{kirchenbauer2023watermark}
J.~Kirchenbauer, J.~Geiping, Y.~Wen, J.~Katz, I.~Miers, and T.~Goldstein.
\newblock A watermark for large language models.
\newblock In \emph{ICML}, 2023.

\bibitem[Mitchell et~al.(2023)]{mitchell2023detectgpt}
E.~Mitchell, Y.~Lee, A.~Khazatsky, C.~D. Manning, and C.~Finn.
\newblock {DetectGPT}: Zero-shot machine-generated text detection using probability curvature.
\newblock In \emph{ICML}, 2023.

\bibitem[Padmakumar and He(2024)]{padmakumar2024does}
V.~Padmakumar and H.~He.
\newblock Does writing with language models reduce content diversity?
\newblock In \emph{ICLR}, 2024.

\bibitem[Shumailov et~al.(2024)]{shumailov2024curse}
I.~Shumailov, Z.~Shumaylov, Y.~Zhao, N.~Papernot, R.~Anderson, and Y.~Gal.
\newblock The curse of recursion: Training on generated data makes models forget.
\newblock \emph{Nature}, 2024.

\bibitem[Solaiman et~al.(2019)]{solaiman2019release}
I.~Solaiman, M.~Brundage, J.~Clark, A.~Askell, A.~Herbert-Voss, J.~Wu, A.~Radford, G.~Krueger, J.~W. Kim, S.~Kreps, et~al.
\newblock Release strategies and the social impacts of language models.
\newblock \emph{arXiv preprint arXiv:1908.09203}, 2019.

\bibitem[Tulchinskii et~al.(2024)]{tulchinskii2024ldscore}
E.~Tulchinskii, K.~Kuznetsov, L.~Kushnareva, D.~Cherniavskii, S.~Nikolenko, E.~Burnaev, S.~Barannikov, and I.~Piontkovskaya.
\newblock Intrinsic dimension estimation for robust detection of {AI}-generated texts.
\newblock In \emph{NeurIPS}, 2024.

\bibitem[Uchendu et~al.(2020)]{uchendu2020authorship}
A.~Uchendu, T.~Le, K.~Shu, and D.~Lee.
\newblock Authorship attribution for neural text generation.
\newblock In \emph{EMNLP}, 2020.

\bibitem[Wenzek et~al.(2020)]{wenzek2020ccnet}
G.~Wenzek, M.-A. Lachaux, A.~Conneau, V.~Chaudhary, F.~Guzm{\'a}n, A.~Joulin, and E.~Grave.
\newblock {CCNet}: Extracting high quality monolingual datasets from web crawl data.
\newblock In \emph{LREC}, 2020.

\bibitem[Wettig et~al.(2024)]{wettig2024qurating}
A.~Wettig, A.~Gao, Z.~Zhong, and D.~Chen.
\newblock {QuRating}: Selecting high-quality data for training language models.
\newblock In \emph{ICML}, 2024.

\bibitem[Xie et~al.(2023)]{xie2023dsir}
S.~M. Xie, S.~Santurkar, T.~Ma, and P.~Liang.
\newblock Data selection for language models via importance resampling.
\newblock In \emph{NeurIPS}, 2023.

\end{thebibliography}

\end{document}
